\chapter{Jump Processes}\label{ch:qm:jump-processes}

On 19 October 1987 the Dow Jones Industrial Average fell 22.6\% in one session. Measured in
the standard deviation of the S\&P~500's daily returns over six decades, 0.98\%, the index's fall
that day was a 21-standard-deviation move. A Gaussian model with that volatility gives such a
day a probability of about $10^{-97}$: waiting for it would take longer than the universe has
existed, many times over. No amount of re-estimating the volatility repairs this; what fails is
the continuity of the paths. A process that moves continuously except at random times, when it
jumps, can give the same day a probability of once in fifty years while keeping the ordinary
days ordinary. This chapter builds the jump processes the series uses: Poisson and compound
Poisson processes, Itô's formula with jumps, and the Lévy processes, characterised by one
formula, the Lévy--Khintchine formula, that later books use to price with (One Quant Book~5,
chapter~13) and to transform (chapter~28 of this book).

\section{Poisson and compound Poisson processes}

\begin{definition}[Poisson process, compound Poisson process]\label{def:qm:jump-processes:poisson}
A \emph{Poisson process}\index{Poisson process} with rate $\lambda > 0$ is a counting process $N_t$ with $N_0 =
0$, independent increments and $N_t - N_s \sim \mathrm{Poisson}(\lambda(t - s))$. A \emph{compound Poisson
process}\index{compound Poisson process} is $J_t = \sum_{i=1}^{N_t}Y_i$, with the jump sizes $Y_i$
independent, identically distributed and independent of $N$.
\end{definition}

The times between events are independent exponentials with mean $1/\lambda$, and the first event
time $\tau_1$ has $\P(\tau_1 > t) = e^{-\lambda t}$: the survival probability of One Quant Book~2,
chapter~23, where the rate is a hazard rate. A \omterm{def:qm:jump-processes:poisson}{Poisson process} is not a \omterm{def:qm:probability-at-speed:martingale}{martingale}, but it
becomes one when its expected growth is removed.

\begin{definition}[Compensated Poisson process]\label{def:qm:jump-processes:compensated}
The \emph{compensated Poisson process}\index{compensated Poisson process} is $\tilde N_t = N_t - \lambda
t$; for a \omterm{def:qm:jump-processes:poisson}{compound Poisson process}, $J_t - \lambda\E[Y]t$.
\end{definition}

Both are \omterm{def:qm:probability-at-speed:martingale}{martingales}, with $\E[\tilde N_t^2] = \lambda t$; the increments are independent and centred,
and the variance of a Poisson variable is its mean. The compensated process is the jump analogue
of \omterm{def:qm:brownian-motion:bm}{Brownian motion}: a \omterm{def:qm:probability-at-speed:martingale}{martingale} with \omterm{def:qm:brownian-motion:qv}{quadratic variation} $[\tilde N]_t = N_t$, the number of jumps
(\cref{fig:qm:jump-processes:paths}, left). Chapter~7 generalises the compensator to rates that
depend on the past.

\section{Itô's formula with jumps}

\begin{definition}[Jump-diffusion, semimartingale]\label{def:qm:jump-processes:jd}
A \emph{jump-diffusion}\index{jump-diffusion} is $X_t = X_0 + \int_0^t\mu_s\,ds + \int_0^t\sigma_s\,dW_s + J_t$,
with $J$ a \omterm{def:qm:jump-processes:poisson}{compound Poisson process} (finitely many jumps on bounded intervals). A
\emph{semimartingale}\index{semimartingale} is a process that is a \omterm{def:qm:ito-calculus:local}{local martingale} plus an \omterm{def:qm:probability-at-speed:filtration}{adapted
process} of finite variation, both càdlàg: the widest class of integrators for which stochastic
integrals and Itô's formula work.
\end{definition}

\begin{theorem}[Itô's formula with jumps]\label{thm:qm:jump-processes:ito}
For a \omterm{def:qm:jump-processes:jd}{jump-diffusion} $X$ and $f \in C^2$,
\[
f(X_t) = f(X_0) + \int_0^tf'(X_{s-})(\mu_s\,ds + \sigma_s\,dW_s) + \tfrac12\int_0^tf^{\prime\prime}(X_{s-})\sigma_s^2\,ds +
\sum_{s \le t}\bigl(f(X_s) - f(X_{s-})\bigr).
\]
\end{theorem}

\begin{proof}
Between two jump times the process is an \omterm{def:qm:ito-calculus:process}{Itô process} and \cref{thm:qm:ito-calculus:ito} applies; at
a jump time $\tau_i$ the change of $f(X)$ is $f(X_{\tau_i}) - f(X_{\tau_i -})$. There are finitely many jumps
on $[0, t]$, so summing the pieces gives the formula.
\end{proof}

The jump term is the exact change, not a Taylor expansion: jumps are not small. With $f = \exp$ and
$S = e^X$, $dS_t/S_{t-} = (\mu_t + \tfrac12\sigma_t^2)\,dt + \sigma_t\,dW_t + (e^{\Delta X_t} - 1)$; compensating the
last term, $S$ is a \omterm{def:qm:probability-at-speed:martingale}{martingale} if and only if $\mu + \tfrac12\sigma^2 + \lambda(\E[e^Y] - 1) = 0$, the drift
correction that every jump model of One Quant Book~5 carries.

\begin{omfigure}
\begin{tikzpicture}
\begin{groupplot}[group style={group size=2 by 1,horizontal sep=1.6cm},
  width=7.6cm,height=4.8cm,tick label style={font=\small},label style={font=\small},
  grid=major,grid style={gray!25},table/col sep=comma]
\nextgroupplot[xlabel={$t$ (years)},xmin=0,xmax=2,ymin=-2,ymax=12,
  legend style={font=\footnotesize,at={(0.03,0.97)},anchor=north west,fill=white,draw=none}]
\addplot[omDef,thick,const plot] table[x=t,y=n]{figdata/methods/06-jump-processes/poisson.csv};
\addplot[omThm,thick] table[x=t,y=comp]{figdata/methods/06-jump-processes/poisson.csv};
\legend{$N_t$,$N_t - 5t$}
\nextgroupplot[xlabel={$t$ (years)},ylabel={log price},xmin=0,xmax=10,ymin=-0.6,ymax=0.3]
\addplot[omDef,thin] table[x=t,y=x]{figdata/methods/06-jump-processes/jd_path.csv};
\addplot[omThm,only marks,mark=o,mark size=2.6pt] table[x=t,y=x]{figdata/methods/06-jump-processes/jd_jumps.csv};
\end{groupplot}
\end{tikzpicture}

\omcaption{Left: a \omterm{def:qm:jump-processes:poisson}{Poisson process} with rate 5 a year and its compensated \omterm{def:qm:probability-at-speed:martingale}{martingale} $N_t - 5t$.
Right: ten years of the chapter's calibrated \omterm{def:qm:jump-processes:jd}{jump-diffusion} (daily steps; jumps at rate 0.5 a year
with mean $-11.1\%$, marked), with the drift that makes the mean daily return zero. Data: the
chapter's tutorial, seeded.}\label{fig:qm:jump-processes:paths}
\end{omfigure}

\section{Lévy processes and the Lévy--Khintchine formula}

\begin{definition}[Lévy process, infinitely divisible distribution]\label{def:qm:jump-processes:levy}
A \emph{Lévy process}\index{Lévy process} is a càdlàg process with $X_0 = 0$ and independent, stationary
increments. A law $\mu$ is an \emph{infinitely divisible distribution}\index{infinitely divisible
distribution} if for every $n$ it is the law of a sum of $n$ independent, identically distributed
variables.
\end{definition}

The law of $X_1$ of a \omterm{def:qm:jump-processes:levy}{Lévy process} is infinitely divisible ($X_1$ is the sum of $n$ increments of length
$1/n$), and conversely every infinitely divisible law is the law of $X_1$ for some \omterm{def:qm:jump-processes:levy}{Lévy process}. \omterm{def:qm:brownian-motion:drift}{Brownian
motion with drift}, the Poisson and compound Poisson processes, and their sums are Lévy processes; the
general one adds infinitely many small jumps.

\begin{definition}[Lévy measure, Lévy triplet, characteristic exponent]\label{def:qm:jump-processes:triplet}
The \emph{Lévy measure}\index{Lévy measure} of a \omterm{def:qm:jump-processes:levy}{Lévy process} is $\nu(A) = \E[\#\{s \le 1 : \Delta X_s \in A\}]$,
the expected number of jumps per unit time with size in $A$ ($0 \notin \bar A$); it satisfies
$\int\min(1, x^2)\,\nu(dx) < \infty$. The \emph{Lévy triplet}\index{Lévy triplet} is $(\sigma^2, \nu, \gamma)$:
Gaussian variance, Lévy measure and drift. The \emph{characteristic exponent}\index{characteristic
exponent} is the function $\psi$ with $\varphi_{X_t}(u) = e^{t\psi(u)}$.
\end{definition}

\begin{theorem}[Lévy--Khintchine]\label{thm:qm:jump-processes:lk}
The \omterm{def:qm:jump-processes:triplet}{characteristic exponent} of a \omterm{def:qm:jump-processes:levy}{Lévy process} with triplet $(\sigma^2, \nu, \gamma)$ is
\[
\psi(u) = \iu\gamma u - \tfrac12\sigma^2u^2 + \int_{\R}\bigl(e^{\iu ux} - 1 - \iu ux\mathbf 1_{\{|x|<1\}}\bigr)\,\nu(dx),
\]
and every triplet with $\sigma^2 \ge 0$ and $\int\min(1, x^2)\,\nu(dx) < \infty$ defines a \omterm{def:qm:jump-processes:levy}{Lévy process}.
Moreover $X$ decomposes into a \omterm{def:qm:brownian-motion:drift}{Brownian motion with drift}, a \omterm{def:qm:jump-processes:poisson}{compound Poisson process} of the jumps
larger than one, and a square-integrable \omterm{def:qm:probability-at-speed:martingale}{martingale} of the compensated small jumps (Lévy--Itô).
\end{theorem}

\admitted

For a \omterm{def:qm:jump-processes:poisson}{compound Poisson process} with rate $\lambda$ and jump law $F$, $\nu = \lambda F$ and the integral is
$\lambda(\varphi_Y(u) - 1)$, up to a drift: the computation is one line of conditioning on $N_t$. The chapter's
\omterm{def:qm:jump-processes:jd}{jump-diffusion}, with Gaussian jumps, has $\psi(u) = \iu\mu u - \tfrac12\sigma^2u^2 + \lambda(e^{\iu u\mu_J - \sigma_J^2u^2/2} -
1)$. Every model of the running project is specified by its exponent (the tutorial's first
listing), and \cref{fig:qm:jump-processes:densities} compares four of them at one day.

\begin{definition}[Subordinator]\label{def:qm:jump-processes:subordinator}
A \emph{subordinator}\index{subordinator} is a \omterm{def:qm:jump-processes:levy}{Lévy process} with nondecreasing paths; running a \omterm{def:qm:brownian-motion:bm}{Brownian
motion} on it as a clock, $X_t = \theta G_t + \sigma W_{G_t}$, gives a \omterm{def:qm:jump-processes:levy}{Lévy process} whose exponent is the
subordinator's Laplace exponent evaluated at $\iu u\theta - \sigma^2u^2/2$.
\end{definition}

The gamma \omterm{def:qm:jump-processes:subordinator}{subordinator} gives the variance gamma process and the inverse Gaussian one the normal
inverse \omterm{def:qm:brownian-motion:bm}{Gaussian process}: trading time that runs fast on busy days and slow on quiet ones. One Quant
Book~5, chapter~13, defines the pricing models built on them. With the same annual variance as the
chapter's \omterm{def:qm:jump-processes:jd}{jump-diffusion}, they spread their tail differently: the \omterm{def:qm:jump-processes:jd}{jump-diffusion} keeps a Gaussian body
and adds a distant bump of crash days, the subordinated models fatten the whole tail smoothly
(\cref{fig:qm:jump-processes:densities}).

\begin{omfigure}
\begin{tikzpicture}
\begin{axis}[width=11cm,height=5.2cm,ymode=log,
  xlabel={one-day log return (\%)},ylabel={density},
  tick label style={font=\small},label style={font=\small},
  xmin=-25,xmax=8,ymin=1e-4,ymax=3,grid=major,grid style={gray!25},unbounded coords=jump,
  legend columns=4,legend style={font=\footnotesize,at={(0.5,-0.3)},anchor=north,draw=none,fill=none,/tikz/every even column/.append style={column sep=6pt}},
  table/col sep=comma]
\addplot[black,thick,dashed] table[x=x,y=gauss]{figdata/methods/06-jump-processes/densities.csv};
\addplot[omThm,thick] table[x=x,y=jd]{figdata/methods/06-jump-processes/densities.csv};
\addplot[omDef,thick] table[x=x,y=vg]{figdata/methods/06-jump-processes/densities.csv};
\addplot[omProp,thick] table[x=x,y=nig]{figdata/methods/06-jump-processes/densities.csv};
\legend{Gaussian,jump-diffusion,variance gamma,normal inverse Gaussian}
\end{axis}
\end{tikzpicture}

\omcaption{Densities of one day's log return, on a log scale, for four Lévy models with the same
variance (a daily standard deviation of 0.98\%): the Gaussian, the chapter's calibrated \omterm{def:qm:jump-processes:jd}{jump-diffusion},
variance gamma ($\theta = -0.1$, $\nu = 0.013$) and normal inverse Gaussian ($\alpha = 58$, $\beta = -11.6$), the last two with one-day excess kurtoses of 8.5 and 12.5. Histograms of
2 million simulated days each, with bins of 0.5 point. Data: the chapter's tutorial,
seeded.}\label{fig:qm:jump-processes:densities}
\end{omfigure}

\section{Cumulants and the term structure of kurtosis}

\begin{definition}[Cumulant]\label{def:qm:jump-processes:cumulant}
The $n$-th \emph{cumulant}\index{cumulant} $\kappa_n$ of a random variable is the $n$-th derivative at
zero of its cumulant generating function $\ln\E[e^{sX}]$; $\kappa_1$ is the mean, $\kappa_2$ the
variance, $\kappa_3/\kappa_2^{3/2}$ the skewness and $\kappa_4/\kappa_2^2$ the excess kurtosis.
\end{definition}

\begin{proposition}[Cumulants of a Lévy process scale with time]\label{prop:qm:jump-processes:scaling}
For a \omterm{def:qm:jump-processes:levy}{Lévy process}, $\kappa_n(X_t) = t\,\kappa_n(X_1)$. Hence the skewness of $X_t$ decays like
$t^{-1/2}$ and its excess kurtosis like $t^{-1}$. For Gaussian jumps, $\kappa_n(X_1)$ is the Gaussian
part plus $\lambda\E[Y^n]$.
\end{proposition}

\begin{proof}
$\ln\E[e^{sX_t}] = t\psi(-\iu s)$ is linear in $t$, and so are its derivatives. For the compound Poisson
part the \omterm{def:qm:jump-processes:cumulant}{cumulant} generating function is $\lambda(\E[e^{sY}] - 1)$, whose $n$-th derivative at zero is
$\lambda\E[Y^n]$.
\end{proof}

A jump model that explains one-day crashes is close to Gaussian at a month: the kurtosis it needs at
one day is diluted by averaging. For the chapter's calibration the excess kurtosis is 76 at one day
and 3.6 at a month (\cref{fig:qm:jump-processes:kurtosis}). Returns measured in markets keep more
kurtosis at a month than this $1/t$ law allows, because their volatility itself clusters: the
independence of increments, not the jumps, is what fails at longer horizons (chapter~18).

\begin{omfigure}
\begin{tikzpicture}
\begin{groupplot}[group style={group size=2 by 1,horizontal sep=1.7cm},
  width=7.6cm,height=4.9cm,tick label style={font=\small},label style={font=\small},
  grid=major,grid style={gray!25},table/col sep=comma]
\nextgroupplot[xlabel={one-day fall (\%)},ylabel={$\log_{10}\P(\text{fall} \ge x)$},xmin=0,xmax=25,ymin=-100,ymax=5,
  legend style={font=\footnotesize,at={(0.03,0.03)},anchor=south west,fill=white,draw=none}]
\addplot[omDef,very thick] table[x=drop,y=normal]{figdata/methods/06-jump-processes/tail.csv};
\addplot[omThm,very thick] table[x=drop,y=jd]{figdata/methods/06-jump-processes/tail.csv};
\addplot[gray,dashed,forget plot] coordinates {(20,-100) (20,5)};
\legend{Gaussian,jump-diffusion}
\nextgroupplot[xlabel={horizon (days)},ylabel={excess kurtosis},xmode=log,ymode=log,xmin=0.8,xmax=320,ymin=0.2,ymax=120,
  legend style={font=\footnotesize,at={(0.97,0.97)},anchor=north east,fill=white,draw=none}]
\addplot[omDef,very thick] table[x=h,y=theory]{figdata/methods/06-jump-processes/kurtosis.csv};
\addplot[omThm,only marks,mark=*,mark size=1.7pt] table[x=h,y=sim]{figdata/methods/06-jump-processes/kurtosis.csv};
\legend{$\kappa_4/(t\,\kappa_2^2)$,simulated}
\end{groupplot}
\end{tikzpicture}

\omcaption{Left: probability that one day's return falls by at least $x$, on a $\log_{10}$ scale, for a
Gaussian with the market's daily standard deviation (0.98\%) and for the chapter's \omterm{def:qm:jump-processes:jd}{jump-diffusion} with the
same variance; at 20\% (dashed) they differ by 88 orders of magnitude. Right: the \omterm{def:qm:jump-processes:jd}{jump-diffusion}'s excess
kurtosis over horizons of 1 to 252 days, from its \omterm{def:qm:jump-processes:cumulant}{cumulants} and from 400\,000 simulated returns per
horizon. Data: the chapter's tutorial, seeded.}\label{fig:qm:jump-processes:kurtosis}
\end{omfigure}

\section{Tutorial: a characteristic function you can check}\label{tut:qm:jump-processes:charfn}

\begin{tutorial}
\textbf{Goal.} Specify the chapter's \omterm{def:qm:jump-processes:jd}{jump-diffusion} by its \omterm{def:qm:jump-processes:triplet}{characteristic exponent}, simulate it, and
check the exponent, the \omterm{def:qm:jump-processes:cumulant}{cumulants} and the kurtosis term structure against the paths.
\textbf{End state:} \cref{fig:qm:jump-processes:paths,fig:qm:jump-processes:kurtosis,fig:qm:jump-processes:charfn}
and the calibration of the weekend problem.
\begin{tutsteps}
\item \textbf{The exponents} of the running project's Lévy models.
\omcode{code/firm/levy/firm_levy.py}{26}{48}{Characteristic exponents: Brownian motion, Gaussian and double-exponential jumps, variance gamma, normal inverse Gaussian.}\label{lst:qm:jump-processes:psi}
\item \textbf{The calibration}: bisect on the mean jump size until a fall of 20\% is a
once-in-fifty-years event, holding the daily variance at $0.98\%^2$.
\omcode{code/methods/06-jump-processes/python/qm_jumps.py}{45}{61}{Calibrating the jump size to a tail frequency at fixed variance.}\label{lst:qm:jump-processes:calibrate}
\item \textbf{Run} \texttt{problem()}, \texttt{charfn\_check()} and \texttt{fig\_jumps.py}; the numerical
\omterm{def:qm:jump-processes:cumulant}{cumulants} of \texttt{firm\_levy.cumulants} agree with the closed forms to $10^{-7}$.
\end{tutsteps}
\textbf{What to change next.} Replace the Gaussian jumps by double-exponential ones with the same mean and
variance and compare the two tails at 20\%; simulate variance gamma by subordination and check its
kurtosis against $3\nu$ at unit time.
\end{tutorial}

\begin{omfigure}
\begin{tikzpicture}
\begin{axis}[width=11cm,height=4.6cm,
  xlabel={$u$},ylabel={$\varphi_{X_t}(u)$},
  tick label style={font=\small},label style={font=\small},
  xmin=0,xmax=40,ymin=-0.1,ymax=1.05,grid=major,grid style={gray!25},
  legend columns=2,legend style={font=\footnotesize,at={(0.5,-0.3)},anchor=north,draw=none,fill=none,/tikz/every even column/.append style={column sep=8pt}},
  table/col sep=comma]
\addplot[omDef,very thick] table[x=u,y=re_th]{figdata/methods/06-jump-processes/charfn.csv};
\addplot[omDef,only marks,mark=*,mark size=1.4pt] table[x=u,y=re_emp]{figdata/methods/06-jump-processes/charfn.csv};
\addplot[omThm,very thick] table[x=u,y=im_th]{figdata/methods/06-jump-processes/charfn.csv};
\addplot[omThm,only marks,mark=*,mark size=1.4pt] table[x=u,y=im_emp]{figdata/methods/06-jump-processes/charfn.csv};
\legend{real part: $e^{t\psi(u)}$,real part: simulated,imaginary part: $e^{t\psi(u)}$,imaginary part: simulated}
\end{axis}
\end{tikzpicture}

\omcaption{\omterm{def:qm:probability-at-speed:charfn}{Characteristic function} of the calibrated \omterm{def:qm:jump-processes:jd}{jump-diffusion}'s one-month return: the
Lévy--Khintchine formula (lines) against the empirical average of $e^{\iu uX_t}$ over 200\,000 simulated
returns (dots). Data: the chapter's tutorial, seeded.}\label{fig:qm:jump-processes:charfn}
\end{omfigure}

\section{Build: the Lévy toolkit}\label{bld:qm:jump-processes:levy}

\begin{build}
\bfield{Purpose} Every jump model of the miniature firm is specified once, by its \omterm{def:qm:jump-processes:triplet}{characteristic
exponent}, and used three ways: to simulate, to compute moments, and to price by transform
(chapter~28).
\bfield{Interface} \texttt{psi\_bm}, \texttt{psi\_merton}, \texttt{psi\_kou}, \texttt{psi\_vg},
\texttt{psi\_nig}; \texttt{cumulants(psi, n\_max)}; \texttt{merton\_cumulants(\dots)};
\texttt{exp\_martingale\_drift(psi)}; \texttt{simulate\_merton}, \texttt{simulate\_vg},
\texttt{simulate\_nig}.
\bfield{Rules} Exponents take complex arguments and follow the series' convention $\varphi_{X_t}(u) =
e^{t\psi(u)}$; simulations are exact on their grid; the \omterm{def:qm:probability-at-speed:martingale}{martingale} drift is computed from the exponent,
never typed.
\bfield{Acceptance tests} \texttt{code/firm/levy/tests/}: numerical \omterm{def:qm:jump-processes:cumulant}{cumulants} equal the closed forms;
empirical \omterm{def:qm:probability-at-speed:charfn}{characteristic functions} of simulated paths match $e^{t\psi(u)}$; variance gamma and normal
inverse Gaussian paths have the moments their exponents give; $e^X$ with the computed drift has constant
mean.
\bfield{Stretch} CGMY and stable exponents; exact simulation of \omterm{def:qm:brownian-motion:passage}{first-passage times} for spectrally
negative processes.
\end{build}

\begin{omsources}
\item Federal Reserve History, ``Stock Market Crash of 1987''.
\item CFA Institute, Enterprising Investor, ``Fact file: S\&P~500 sigma events'', 2012.
\item R.~C.~Merton, ``Option pricing when underlying stock returns are discontinuous'',
\emph{Journal of Financial Economics} 3, 1976.
\item S.~G.~Kou, ``A jump-diffusion model for option pricing'', \emph{Management Science} 48, 2002.
\item R.~Cont and P.~Tankov, \emph{Financial Modelling with Jump Processes}, Chapman \& Hall/CRC, 2004.
\end{omsources}

\section{Exercises}

\begin{exercise}[$\star$]\label{exo:qm:jump-processes:1}
Events arrive as a \omterm{def:qm:jump-processes:poisson}{Poisson process} at 5 a year. What is the probability of none in a month, and the
expected wait for the next one?
\end{exercise}

\begin{exercise}[$\star$]\label{exo:qm:jump-processes:2}
Jumps arrive at 3 a year with sizes $\mathcal N(-2\%, 4\%^2)$. What are the mean and standard deviation of their
sum over a year?
\end{exercise}

\begin{exercise}[$\star$]\label{exo:qm:jump-processes:3}
Show that $N_t - \lambda t$ is a \omterm{def:qm:probability-at-speed:martingale}{martingale} and that $\E[(N_t - \lambda t)^2] = \lambda t$.
\end{exercise}

\begin{exercise}[$\star\star$]\label{exo:qm:jump-processes:4}
For $X = \sigma W + J$ with $J$ compound Poisson, write $d(e^{X_t})$ with \cref{thm:qm:jump-processes:ito}, and
find the drift to add to $X$ so that $e^X$ is a \omterm{def:qm:probability-at-speed:martingale}{martingale}.
\end{exercise}

\begin{exercise}[$\star\star$]\label{exo:qm:jump-processes:5}
Jumps arrive at rate 2 with exponential sizes of mean $1/50$. Write the \omterm{def:qm:jump-processes:triplet}{characteristic exponent} of the
\omterm{def:qm:jump-processes:poisson}{compound Poisson process} and its mean and variance at $t = 1$.
\end{exercise}

\begin{exercise}[$\star\star$]\label{exo:qm:jump-processes:6}
A Lévy model has an excess kurtosis of 76.4 at one day. What does it have at a week of five days and at
a year of 252?
\end{exercise}

\begin{exercise}[$\star\star\star$]\label{exo:qm:jump-processes:7}
\emph{Coding.} With \texttt{cumulants}, compute the first four \omterm{def:qm:jump-processes:cumulant}{cumulants} of the variance gamma process
with $\theta = -0.1$, $\sigma = 0.2$, $\nu = 0.3$ at unit time, and check them against $\kappa_1 = \theta$, $\kappa_2 = \sigma^2 +
\nu\theta^2$, $\kappa_3 = 2\theta^3\nu^2 + 3\sigma^2\theta\nu$, $\kappa_4 = 3\sigma^4\nu + 6\theta^4\nu^3 + 12\sigma^2\theta^2\nu^2$.
\end{exercise}

\begin{exercise}[$\star\star\star$]\label{exo:qm:jump-processes:8}
\emph{Find the flaw.} ``Monthly returns of our index have an excess kurtosis of 1.2, close to Gaussian,
so jumps do not matter for our one-day risk.''
\end{exercise}

\section{Problem: Twenty Standard Deviations}

\begin{problem}[{Weekend problem --- a crash that a diffusion cannot produce}]\label{pb:qm:jump-processes:1}
Daily returns of an index have a standard deviation of 0.98\%. A risk manager wants a model in which a
one-day fall of 20\% or more happens once in fifty years of 252 days, and ordinary days keep that
variance. The model is a \omterm{def:qm:jump-processes:jd}{jump-diffusion}: jumps at rate $\lambda = 0.5$ a year with sizes $\mathcal N(\mu_J,
0.05^2)$, and a diffusion whose volatility makes up the rest of the variance.

\medskip\noindent\textbf{Part I --- The day.}
\begin{enumerate}
\item How large a fall is 20.98 standard deviations?
\item What probability does a Gaussian with that standard deviation give it? (Give $\log_{10}$.)
\item And a fall of exactly 20\%?
\item What annual volatility does 0.98\% a day make?
\item Why can re-estimating the volatility not rescue the Gaussian model?
\end{enumerate}

\medskip\noindent\textbf{Part II --- The \omterm{def:qm:jump-processes:jd}{jump-diffusion}.}
\begin{enumerate}[resume]
\item What mean jump size $\mu_J$ makes a 20\% fall a once-in-fifty-years event?
\item What diffusion volatility keeps the daily variance at $0.98\%^2$?
\item What share of the variance do the jumps carry?
\item What is the probability of a jump on a given day?
\item How often does the model produce a fall of 10\% or more?
\end{enumerate}

\medskip\noindent\textbf{Part III --- \omterm{def:qm:jump-processes:cumulant}{Cumulants}.}
\begin{enumerate}[resume]
\item What are the one-day skewness and excess kurtosis?
\item What is the excess kurtosis at a month of 21 days?
\item Why does it fall like $1/t$?
\item What drift correction makes $e^X$ a \omterm{def:qm:probability-at-speed:martingale}{martingale}?
\item Check the kurtosis at five days against simulation.
\end{enumerate}

\medskip\noindent\textbf{Part IV --- Judgement.}
\begin{enumerate}[resume]
\item Is ``once in fifty years'' an estimate or a choice?
\item What would the model miss that the data show at a month?
\item Which options does the calibration matter most for?
\item State the \emph{named result}: the size of the 1987 day under a diffusion, and the jump that makes it rare but possible.
\item In one sentence: what does a jump process change that a volatility cannot?
\end{enumerate}
\end{problem}

\section{Interview questions}

\begin{interviewq}[$\star$ \iqroles{trader, risk}]\label{iq:qm:jump-processes:1}
A model says yesterday's move was a 20-standard-deviation event. What do you conclude?
\end{interviewq}

\begin{interviewq}[$\star$ \iqroles{researcher}]\label{iq:qm:jump-processes:2}
Trades arrive as a \omterm{def:qm:jump-processes:poisson}{Poisson process} at 2 a second. What is the probability of at least two in a second, and
the expected time to the next one?
\end{interviewq}

\begin{interviewq}[$\star\star$ \iqroles{bank, trader}]\label{iq:qm:jump-processes:3}
Why can a delta hedge not remove the risk of a jump?
\end{interviewq}

\begin{interviewq}[$\star\star$ \iqroles{researcher}]\label{iq:qm:jump-processes:4}
State the Lévy--Khintchine formula and say what each part of the triplet means.
\end{interviewq}

\begin{interviewq}[$\star\star$ \iqroles{researcher, risk}]\label{iq:qm:jump-processes:5}
How does the kurtosis of returns change with the horizon in a Lévy model, and in the data?
\end{interviewq}

\begin{interviewq}[$\star\star\star$ \iqroles{bank, researcher}]\label{iq:qm:jump-processes:6}
In a \omterm{def:qm:jump-processes:jd}{jump-diffusion} for a stock under the \omterm{def:qm:girsanov-and-changes-of-numeraire:emm}{risk-neutral measure}, what must the drift be, and why?
\end{interviewq}
